\subsection{同态 F 与移位 V}

设 A 为一个环。在本段其余部分中，我们分别用 + 和 × 表示 W(A) 中的加法和乘法运算。我们还将 0 记作 0 \(_{A}\)，将 1 记作 \(1_{A}\)。我们按下列公式定义从 W(A) 到自身的两个映射 \(^{1}\) \(F_{A}\) 和 \(V_{A}\)（也简记为 F 和 V）。

\begin{enumerate}
\def\labelenumi{(\arabic{enumi})}
\setcounter{enumi}{22}
\tightlist
\item
\end{enumerate}

\[
\mathrm{F} _ {\mathbf {A}} (\boldsymbol {a}) = \big (\mathrm{F} _ {n} (a _ {0},..., a _ {n + 1}) \big) _ {n \in \mathbb {N}},\tag{24}
\]

\[
\mathrm{V} _ {\mathrm{A}} (\boldsymbol {a}) = (0, a _ {0}, a _ {1}, \dots)
\]

\((\text{其中 } \boldsymbol{a} = (a_n)_{n \in \mathbb{N}} \text{ 属于 } \mathrm{W(A)})\)。映射 \(V_A\) 称为移位。公式

\[
\Phi_ {n} \left(\mathrm{F} _ {0} (\boldsymbol {a}), \dots , \mathrm{F} _ {n} (\boldsymbol {a})\right) = \Phi_ {n + 1} \left(a _ {0}, \dots , a _ {n + 1}\right) \quad (n \in \mathbf {N})\tag{25}
\]

由 (15) 立即得到。它也可写成如下形式：

\[
\Phi_ {\mathrm{A}} \circ \mathrm{F} _ {\mathrm{A}} = f _ {\mathrm{A}} \circ \Phi_ {\mathrm{A}}.\tag{26}
\]

公式

\[
\Phi_ {\mathrm{A}} \circ \mathrm{V} _ {\mathrm{A}} = v _ {\mathrm{A}} \circ \Phi_ {\mathrm{A}}\tag{27}
\]

由关系 (3) 得到。

设 \(\rho : \mathbf{B} \to \mathbf{A}\) 为一个环同态。关系式

\begin{enumerate}
\def\labelenumi{(\arabic{enumi})}
\setcounter{enumi}{27}
\tightlist
\item
\end{enumerate}

\[
\mathbf {W} (\rho) \circ \mathrm{F} _ {\mathbf {B}} = \mathrm{F} _ {\mathbf {A}} \circ \mathbf {W} (\rho)\tag{29}
\]

\[
\mathrm{W} (\rho) \circ \mathrm{V} _ {\mathrm{B}} = \mathrm{V} _ {\mathrm{A}} \circ \mathrm{W} (\rho)
\]

由定义立即得到。

命题 3．——设 A 为一个环。

\begin{enumerate}
\def\labelenumi{\alph{enumi})}
\item
  映射 \(F_{A}\) 是环 W(A) 的自同态。
\item
  映射 \(V_{A}\) 是环 \(\mathbf{W}(\mathbf{A})\) 的底层加法群的自同态。
\item
  对 W(A) 中的每个 a，都有 \(F_{A}(V_{A}(a)) = p.a\)（即 W(A) 中 p 个等于 a 的项之和）。
\item
  对 W(A) 中的任意 \(a\) 和 \(b\)，都有
\end{enumerate}

\begin{enumerate}
\def\labelenumi{(\arabic{enumi})}
\setcounter{enumi}{29}
\tightlist
\item
\end{enumerate}

\[
\mathrm{V} _ {\mathrm{A}} (a \times \mathrm{F} _ {\mathrm{A}} (b)) = \mathrm{V} _ {\mathrm{A}} (a) \times b\tag{31}
\]

\[
\mathrm{V} _ {\mathrm{A}} (a) \times \mathrm{V} _ {\mathrm{A}} (b) = p. \mathrm{V} _ {\mathrm{A}} (a \times b)
\]

（即 W(A) 中 p 个等于 \(V_A(a \times b)\) 的项之和）。

\begin{enumerate}
\def\labelenumi{\alph{enumi})}
\setcounter{enumi}{4}
\tightlist
\item
  置 \(\mu = V_{A}(1) = (0, 1, 0, \ldots)\)。对 W(A) 中的每个 \(b\)，都有
\end{enumerate}

\[
\mathrm{V} _ {\mathrm{A}} \left(\mathrm{F} _ {\mathrm{A}} (\boldsymbol {b})\right) = \mu \times \boldsymbol {b}.\tag{32}
\]

\begin{enumerate}
\def\labelenumi{\alph{enumi})}
\setcounter{enumi}{5}
\tightlist
\item
  对 W(A) 的每个元素 \(\mathbf{a}\)，记 W(A) 中 p 个等于 a 的元素的乘积为 \(\mathbf{a}^{*p}\)。则有
\end{enumerate}

\[
\mathrm{F} _ {\mathbf {A}} (\boldsymbol {a}) \equiv \boldsymbol {a} ^ {* p} \bmod p. \mathrm{W} (\mathbf {A}) \quad (\text { ideal of } \mathrm{W} (\mathbf {A}) \text { generated by } p. \mathbf {1}).\tag{33}
\]

设 \(\rho: \mathbf{B} \to \mathbf{A}\) 为满足第 4 节引理 3 条件的环同态。于是 \(\mathbf{W}(\rho): \mathbf{W}(\mathbf{B}) \to \mathbf{W}(\mathbf{A})\) 是满射环同态，而 \(\Phi_{\mathbf{B}}: \mathbf{W}(\mathbf{B}) \to \mathbf{B}^{\mathbf{N}}\) 是单射环同态。此外，\(f_{\mathbf{B}}: \mathbf{B}^{\mathbf{N}} \to \mathbf{B}^{\mathbf{N}}\) 是环同态。由公式 (26) 和 (28)，有

\[
\Phi_ {\mathrm{B}} \circ \mathrm{F} _ {\mathrm{B}} = f _ {\mathrm{B}} \circ \Phi_ {\mathrm{B}}, \quad \mathrm{W} (\rho) \circ \mathrm{F} _ {\mathrm{B}} = \mathrm{F} _ {\mathrm{A}} \circ \mathrm{W} (\rho),
\]

因此断言 a) 立即成立。断言 b) 同理，由公式 (27)、(29) 以及 \(v_{B}\) 是 \(B^{N}\) 的底层加法群的自同态这一事实得到。

设 \(\pmb{a}\) 是 W(A) 的一个元素，并取 W(B) 中的元素 \(\pmb{x}\)，使其在 W(ρ) 下的像为 \(\pmb{a}\)。置 \(\xi = \Phi_{\mathrm{B}}(\pmb{x})\)。由 \(f_{\mathrm{B}}\) 和 \(v_{\mathrm{B}}\) 的定义立即得到 \(f_{\mathrm{B}}(v_{\mathrm{B}}(\xi)) = p.\xi\)（即 B \(^{\mathbf{N}}\) 中 \(p\) 个等于 \(\xi\) 的项之和）。根据公式 (26)、(27)（其中以 B 代替 A），W(B) 中的元素 F \(_{\mathrm{B}}\) (V \(_{\mathrm{B}}\) ( \(\pmb{x}\) )) 和 \(p.\pmb{x}\) 具有相同的像 \(p.\xi\)（经由单射 \(\Phi_{\mathrm{B}}\)），因而相等。于是，公式 F \(_{\mathrm{A}}\) (V \(_{\mathrm{A}}\) ( \(\pmb{a}\) )) = \(p.\pmb{a}\) 由关系 (28)、(29) 得到。这就证明了 \(c\)。

同理，公式 (30) 的证明归结为证明关系式

\[
v _ {\mathrm{B}} (\xi f _ {\mathrm{B}} (\eta)) = v _ {\mathrm{B}} (\xi) \eta
\]

其中 \(\xi\)、\(\eta\) 属于 B \(^{N}\)。这由等式

\[
\begin{array}{l} \xi f _ {\mathrm{B}} (\eta) = (\xi_ {0} \eta_ {1}, \xi_ {1} \eta_ {2}, \dots) \\ v _ {\mathrm{B}} (\xi) \eta = (0, p \xi_ {0} \eta_ {1}, p \xi_ {1} \eta_ {2}, \dots). \end{array}
\]

考虑到 (b)、(c)，公式 (31) 由公式 (30) 得到，其中 \(\boldsymbol{b}\) 被 \(\mathrm{V}_{\mathbf{A}}(\boldsymbol{b})\) 所代替。公式 (32) 是公式 (30) 在 \(\boldsymbol{a} = \boldsymbol{1}\) 时的特例。

类似地，公式 (33) 的证明归结为证明关系式

\[
f _ {\mathrm{B}} (\xi) \equiv \xi^ {p} \bmod p. \Phi_ {\mathrm{B}} (\mathrm{B} ^ {\mathbf {N}}),
\]

其中 \(\xi^{p}\) 表示 \(\mathbf{B}^{\mathbf{N}}\) 中 \(p\) 个等于 \(\xi\) 的元素的乘积。根据第 2 节命题 2 的 c)，这等价于：对每个 \(n \geqslant 0\)，都有

\[
\sigma \left(\xi_ {n + 1} - \xi_ {n} ^ {p}\right) \equiv \xi_ {n + 2} - \xi_ {n + 1} ^ {p} \bmod p ^ {n + 2} \mathbf {B}.
\]

现在，对每个 \(n \geqslant 0\)，由上引文有

\[
\sigma (\xi_ {n}) \equiv \xi_ {n + 1} \bmod p ^ {n + 1} \mathbf {B}
\]

因为 \(\xi = \Phi_{\mathbf{B}}(\boldsymbol{x})\)；由此借助第 1 节引理 1 可推出

\[
\sigma (\xi_ {n}) ^ {p} \equiv \xi_ {n + 1} ^ {p} \bmod p ^ {n + 2} \mathbf {B}.
\]

这就证明了所需的关系式。

注记．——关于在普适 Witt 环情形下与映射 F、V 类似的映射的定义，见第 52 页及以下的习题 36、37、38。

